%% frontmatter.tex --- inherited front matter followed by current
%% administrative/navigation guidance.  The inherited source title page is
%% retained here after the book's governing publication title page; neither
%% the title attribution nor the later guidance substitutes for, abridges, or
%% limits an inherited source presentation.

% Navigation labels are supplied outside the byte-exact artifacts.  The hooks
% execute at the inherited heading commands and do not rewrite source tokens.
\AddToHookNext{env/titlepage/begin}{%
  \phantomsection\label{sec:inherited-title-page}}
\begin{inheritedrestorationscope}{SRC-00171-00197}

\hypersetup{pageanchor=false}
\begin{titlepage}
  \centering
  \vspace*{2.2cm}
  {\Huge\bfseries Formal Logic\par}
  \vspace{0.6cm}
  {\LARGE Language, Semantics, Deduction, and Metatheory\par}
  \vspace{1.2cm}
  {\Large A Graduate-Level Reference\par}
  \vfill
  {\large Version 1.0\par}
  \vspace{0.2cm}
  {\large 9 August 2026\par}
  \vfill
  \begin{minipage}{0.82\textwidth}
  \small
  This document fixes one technical term for each concept and records the scope
  of every notation before use.  Definitions precede use.
  Every calculus is identified before the derivability symbol is used, and
  every semantic claim states the semantic data over which it is made.
  \end{minipage}
\end{titlepage}

\frontmatter
\hypersetup{pageanchor=true}


\end{inheritedrestorationscope}


\AddToHookNext{cmd/chapter/before}{%
  \phantomsection\label{sec:inherited-preface}}
\begin{inheritedrestorationscope}{SRC-00198-00234}
\chapter*{Preface}
\addcontentsline{toc}{chapter}{Preface}

Formal logic studies formal languages, interpretations of those languages,
and rule-governed derivations.  Correct work in formal logic therefore requires
more than symbol manipulation.  It requires exact control of the level of
language, the grammar of expressions, variable binding, interpretation,
assumptions, inference rules, and the metatheoretic hypotheses under which a
result is stated.

This reference is organized at graduate level rather than as a catalogue of
disconnected systems.  The dependency order is deliberate.  Syntax is
defined before semantics; semantics is defined before soundness; a calculus is
defined before derivability; classical propositional logic is established
before first-order logic; classical first-order logic is established before
model theory, computability, incompleteness, proof theory, and nonclassical
logic.  The organization follows the core treatments in
\cite{Enderton2001,EbbinghausFlumThomas1994,OpenLogic2026,Marker2024}, while the
historical and advanced results are tied to their original sources where
practical.

The document has four purposes:
\begin{enumerate}
  \item to teach the formal-logical content of the subject;
  \item to teach the construction and checking of derivations, proofs,
        countervaluations, structures, and countermodels of sentence
        arguments;
  \item to teach the writing and typesetting of formal logic; and
  \item to provide a dependency-ordered reference for graduate study, teaching,
        and preparation for qualifying examinations.
\end{enumerate}

The reference is self-contained at the level of notation and foundational
definitions.  Advanced model theory and proof theory are presented to the
point required for a broad graduate-level reference, but a research specialization
still requires the monographs cited in the relevant chapters.


\end{inheritedrestorationscope}


\AddToHookNext{cmd/section/before}{%
  \phantomsection\label{sec:inherited-terminology-policy}}
\begin{inheritedrestorationscope}{SRC-00235-00293}
\section*{Terminology and notation policy}

Literature in formal logic does not impose one universal notation.  Pretending
otherwise would be inaccurate.  This document instead fixes a registered
global convention.  A subsystem-specific departure is permitted only when its
scope and standing are declared before use, and no registered notation is used
for a second role within that scope.

Whenever this document introduces a new technical name, five conditions are
met at the same location:
\begin{enumerate}
  \item \emph{Necessity}: an existing term cannot express the required
        distinction without ambiguity;
  \item \emph{Definition}: an exact eliminable definition or an explicit axiom
        is supplied;
  \item \emph{Standing}: the grammatical and formal standing of the
        name is stated;
  \item \emph{Nonredundancy}: the distinction from previously defined terms is
        explained; and
  \item \emph{Use}: later definitions or theorems require the distinction.
\end{enumerate}
Until those conditions are met, the text describes the required object or
relation without assigning it a technical name.

Throughout this document, the word \emph{iff} abbreviates the
metalinguistic biconditional ``if and only if''; it is the registered
connective of definitions and metatheoretic equivalences.

The following distinctions are permanent:
\begin{itemize}
  \item A \emph{formula} is a grammatical object-language expression.  A
        \emph{sentence} is a formula with no free variable occurrence.
  \item A \emph{structure} supplies a domain and interpretations for the
        nonlogical symbols of a signature.  A \emph{model} is a structure that
        satisfies a specified set of sentences.
  \item A \emph{valuation} assigns truth values to propositional variables.  A
        \emph{variable assignment} assigns domain elements to first-order
        variables.
  \item \emph{Satisfaction} is the relation between a formula and the complete
        semantic data permitted by the active semantics.  \emph{Truth} is
        satisfaction relative to specified such data.  \emph{Validity} is
        satisfaction for every permitted tuple of semantic data.
        \emph{Semantic consequence} is preservation of satisfaction from a set
        of assumptions to a conclusion.  For a first-order sentence, assignment
        independence later permits the specialized phrase ``truth in a
        structure.''
  \item A \emph{derivation} is a finite rule-governed object in a specified
        calculus.  A \emph{proof} is a derivation with no undischarged
        assumptions other than the stated premises.
  \item \emph{Soundness} proceeds from derivability to semantic consequence.
        \emph{Completeness} proceeds from semantic consequence to derivability.
  \item \emph{Consistency} excludes derivability of a contradiction.
        A theory is \emph{deductively complete} iff every sentence or its
        negation is derivable from the theory.  A theory is
        \emph{semantically complete} iff it semantically entails every sentence
        or its negation.  For the fixed sound and complete first-order calculus,
        these two properties are equivalent; neither is consistency.
\end{itemize}


\end{inheritedrestorationscope}


\AddToHookNext{cmd/section/before}{%
  \phantomsection\label{sec:inherited-evidence-policy}}
\begin{inheritedrestorationscope}{SRC-00294-00320}
\section*{Evidence and verification}

The status of a formal result must be reported at the level actually checked.
These levels are not interchangeable:
\begin{enumerate}
  \item a displayed calculation checks only its displayed instances;
  \item a truth table checks all valuations of the finitely many variables it
        contains;
  \item a derivation checks derivability in the stated calculus;
  \item a proof establishes the stated result under its stated
        hypotheses;
  \item a proof-assistant kernel checks a formal proof relative to the kernel,
        imported declarations, and admitted axioms;
  \item successful \LaTeX{} compilation checks typesetting syntax, not logical
        correctness; and
  \item editorial or formal-logical review is independent evidence and must not be
        inferred from any preceding level.
\end{enumerate}

When this document states a result status explicitly, it uses exactly one of
three terms.  \emph{Complete proof} means that all substantive formal-logical
steps are supplied in this document.  \emph{Proof sketch} identifies the
construction and states which substantive verification is omitted.
\emph{Cited theorem} states exact hypotheses and a conclusion but refers to a
named source for the proof.  A dependency used within a complete proof is
stated separately as a cited theorem; it is not hidden in a qualified status.


\end{inheritedrestorationscope}


\AddToHookNext{cmd/chapter/before}{%
  \phantomsection\label{sec:inherited-graduate-scope}}
\newcounter{hbrestoredfrontmattersection}
\setcounter{hbrestoredfrontmattersection}{0}
\AddToHook{cmd/section/before}[hb-restored-frontmatter-section-labels]{%
  \stepcounter{hbrestoredfrontmattersection}%
  \ifcase\value{hbrestoredfrontmattersection}\or
    \phantomsection\label{sec:inherited-prerequisites}%
  \or
    \phantomsection\label{sec:inherited-learning-outcomes}%
  \fi}
\begin{inheritedrestorationscope}{SRC-00321-00362}
\chapter*{Graduate-Level Scope}
\addcontentsline{toc}{chapter}{Graduate-Level Scope}

\section*{Prerequisites}

The reader should be able to work with sets, functions, relations, finite
strings, countable sets, induction on natural numbers, and ordinary rigorous
proofs.  Prior exposure to elementary logic is useful but not
required.  Chapters that use cardinal arithmetic, algebra, or topology state
the additional facts they require.

\section*{Learning outcomes}

A reader completing the reference should be able to:
\begin{enumerate}
  \item define a formal language by an inductive grammar and prove results by
        structural induction;
  \item define capture-avoiding substitution and prove a substitution lemma;
  \item prove soundness for a natural-deduction or sequent calculus;
  \item prove completeness for classical propositional logic and explain a
        Henkin proof of completeness for classical first-order logic;
  \item derive compactness and the downward and upward L\"owenheim--Skolem
        theorems with their cardinal hypotheses;
  \item use back-and-forth arguments, elementary chains, and ultraproducts;
  \item prove undecidability by a many-one reduction;
  \item state and apply the diagonal lemma and derive the first incompleteness
        theorem, the second incompleteness theorem, and undefinability of truth
        with exact hypotheses;
  \item explain cut elimination, normalization, and the subformula property;
  \item give relational semantics for intuitionistic and normal modal logic and
        prove representative soundness results;
  \item distinguish full semantics from Henkin semantics for second-order logic;
        and
  \item audit a written or mechanized proof for hidden assumptions, variable
        capture, rule misuse, and shifts between object language and
        metalanguage.
\end{enumerate}

\tableofcontents

\mainmatter


\end{inheritedrestorationscope}

\RemoveFromHook{cmd/section/before}[hb-restored-frontmatter-section-labels]

\chapter*{How to Use This Handbook}
\addcontentsline{toc}{chapter}{How to Use This Handbook}
\markboth{HOW TO USE THIS HANDBOOK}{HOW TO USE THIS HANDBOOK}

\section*{Entries}

The handbook contains numbered lookup entries---definitions, conventions,
methods, examples, warnings, proof-boundary notices, theorems, lemmas, and
corollaries---alongside the inherited presentations.  Inherited proofs, rules,
side conditions, examples, qualifications, and attributions remain at their
source-mapped locations.  Corrections are incorporated into the definitions,
statements, and proofs where the reader uses them; the research notes record
the original readings and the reasons for each correction.
The lookup layer supplies additional navigation.
Every theorem-class entry states its exact
hypotheses---including finiteness, countability, consistency, effectiveness,
and choice-principle assumptions---and remains intelligible when
encountered by page lookup.

\section*{Proof-status taglines}
\index{status tagline}

Every editorial theorem, lemma, and corollary carries a status tagline of
exactly one of three forms.  These tags annotate the entry; they do not replace
or recategorize an inherited result presentation.

\begin{description}
  \item[proved (Annex $P n.k$).] A complete proof of the stated result is
        item $P n.k$ of the proof annex.  All substantive formal-logical
        steps are supplied there; dependencies used inside the proof are
        stated separately as cited results.
  \item[sketched (Annex $P n.k$).] The annex item identifies the
        construction and opens with a note enumerating exactly which
        substantive verifications are omitted.  A sketch is not a proof;
        results resting on a sketch are themselves at most sketched.
  \item[cited (source).] The entry states exact hypotheses and a
        conclusion, and the cited source carries the proof.  Cited results
        have no annex item.
\end{description}

These forms implement the inherited evidence policy above.  A status is never
silently upgraded, and a tagline reports the level checked in this work,
not a stronger level attainable elsewhere.

\section*{The proof annex}
\index{proof annex}

Appendix~P records additional proof-audit presentations in chapters numbered
parallel to Parts~1--12.  It does not relocate, retire, or replace an inherited
proof.  Each annex item names exactly one handbook result, and that entry's
status tagline links back to the item; the correspondence is one-to-one and
checked mechanically.  Sketch items open with an explicit \emph{Omissions}
note.  Commentary attached to a cited result is labeled as commentary and is
not an annex item.

\section*{Calculi and the symbol registry}

The calculus registry provides one navigation sheet per calculus, each sheet
recording its language, sequent shape, negation treatment, and equality
treatment.  The sheets link to their controlling sources and dependencies;
they do not replace or normalize inherited rule displays.
Likewise, the symbol registry presented in Part~1 records each current
handbook glyph's macro, role, scope, precedence, and first use without
replacing any inherited notation presentation.  Every attributed source
presentation constitutes its own notation scope and keeps its published
syntax.  A newly added editorial glyph carries one role within its declared
scope; a scoped departure is declared before use.

\chapter*{Editorial Conventions and Evidence Policy}
\addcontentsline{toc}{chapter}{Editorial Conventions and Evidence Policy}
\markboth{EDITORIAL CONVENTIONS}{EDITORIAL CONVENTIONS}

\section*{Naming policy}

The inherited terminology and notation policy in the restored front matter
remains the source presentation.  The following editorial rule governs only new
handbook terminology added outside the inherited source: a subsystem-specific
departure is permitted only when its scope and standing are declared before
use, and no registered notation is used for a second role within that scope.

Whenever the handbook introduces a new technical name, five conditions are
met at the same location:
\begin{enumerate}
  \item \emph{Necessity}: an existing term cannot express the required
        distinction without ambiguity;
  \item \emph{Definition}: an exact eliminable definition or an explicit
        axiom is supplied;
  \item \emph{Standing}: the grammatical and formal standing of the name is
        stated;
  \item \emph{Nonredundancy}: the distinction from previously defined terms
        is explained; and
  \item \emph{Use}: later definitions or theorems require the distinction.
\end{enumerate}
Without all five conditions, the text states the required object, relation, or
statement without assigning it a technical name.
